% Generated from the closed Phase 4 receipts by render_phase4.py.
% Document rendering only; no model or scientific execution.
\section{Original O service and matched digital apparatus}
\subsection{Unchanged external task and service predicate}
All three O0 episodes use the original 21,600 s A0 supply schedule and zero
initial capacitor voltage. The photovoltaic area is 15 cm$^2$ and capacitance
is fixed at 15 mF. The six light intervals are 200, 0, 50, 200, 0 and 50 lux,
at boundaries 0, 120, 150, 180, 300, 330 and 360 minutes. Their assumed
electrical port powers are respectively 90, 0, 20, 90, 0 and 20 $\mu$W.
The 3500 K light specification is metadata; lux is never converted to watts.
Calibrated spectral irradiance and photovoltaic IV data are unavailable.

External RMS vibration is 0.08 g during 70--75, 130--135, 190--200 and
280--282 minutes and 0.005 g otherwise. Start requires continuously greater
than 0.05 g for 3 s; stop requires continuously below 0.01 g for 60 s.
The corresponding threshold equality resets confirmation. The event time is
the confirmation instant. Every start and stop must reach the receiver within
10 s and have a timestamp error at most 1 s. Received active-state messages
must cover each confirmed active interval with gaps at most 120 s. Darkness
is included; idle false messages fail. A correct final state cannot cancel
an earlier missing event or keepalive~\cite{case}.

The four confirmed active intervals end at 4560, 8160, 12060 and 16980 s.
They are service intervals, not a daily occupancy summary. In particular, the
130--135 minute activity window occurs within the first dark supply interval.
The event predicate neither replaces event history with an hourly state nor
accepts a retained NVM value as a newly observed event.

\subsection{Inherited energy law and proposed RMS costs}
The apparatus retains the A/R assumed source/cold-start/PMIC/capacitor and
digital receiver parameters. Harvest efficiency is 0.8, with 0.38 V and
3 $\mu$W cold-start gates. Ready/off voltages are 2.7/2.3 V; maximum capacitor
voltage is 3.3 V, rating 5.5 V and ESR 0.3 ohm. The managed output is 1.8 V
with 0.3 V dropout. Off/sleep power is 1/5 $\mu$W; boot and storage-side TX
costs are 100 and 130 $\mu$J. Transmission uses the existing ideal atomic
wire/CRC and receiver abstraction, not a qualified RF link.

The O addition declares 0.5 $\mu$W powered standby, 1 $\mu$J per wake,
20 $\mu$J per successful RMS sample and 5 $\mu$J per committed write.
The 20 $\mu$J value borrows the numerical A/R temperature-sample cost; it does
not qualify a vibration sensor. A write is charged when a record is committed,
not at every poll. NVM preserves committed state and the pending queue but
provides no sensing energy. Powered samples use a declared zero-order hold;
failed samples and brownouts reset continuous-confirmation timers.
Energy integration is piecewise and observations are compressed into sampling
runs. These are explicit digital assumptions with physical qualification
UNKNOWN~\cite{case}.

\subsection{Candidate, controls and public observations}
All three episodes share the same source, initial state, components, link,
service predicate and hardware. The legal candidate contains only typed
capacitance, sample period, thresholds, cache, retry/backoff and program
choices. The fixed periodic control takes one RMS snapshot every 60 s and
uploads it without an event confirmation state machine. The rich-state control
uses a one-second sampled start/stop state machine, 120 s active keepalive,
NVM queue, one-second retry and zero TX backoff. Each fixed control is its
episode's sole legal design; the candidate may check at most two distinct
designs under the same registered total allowance.

Every initial O0 episode has a ceiling of 24 provider requests, 40 tool calls,
1200 s and two distinct design exposures. Participants must actually inspect,
check, judge and submit. Source records carry hashes and unit maps. The common
public-response projection replaces every array longer than twelve entries
with its full count, canonical full-value hash, first six and last six entries.
The complete raw digital response remains host-sealed, and public feedback
contains no reference truth or host service label. Thus an evaluator's use of
the full response is separate from what the participant could inspect~\cite{protocol}.
All three episode receipts report the \path{deepseek-flash} provider alias;
the protocol leaves pinned weight identity UNKNOWN.

\section{Frozen procedures, selection and final information}
H0 is the unchanged base design procedure: read the fixed task and available
sources, choose a legal D, actually check it, then submit a source-bound design
and judgment while retaining failures and UNKNOWN. G0 is the fixed parent
improver. A single actual proposal episode forms G1 from the frozen parent
material and retained public history. These are changes to instructional
procedures; no model weights, scientific semantics or device parameters are
trained by the G-formation step.

G1's recorded proposal calls for source-parameterized H text, capability and
budget checks before spending an attempt, separation of achieved delivery from
post-terminal decay, and retrospective rehearsal of quantitative predictions
on their own source episode. It distinguishes arithmetic/transcription errors
from substantive design failures and retains no-update as an option. This
describes the proposed procedure; it does not establish that every consumer
follows it or that any rule is a novel or necessary mechanism~\cite{parents}.

Three lineages each generate two H proposals under G0 and two under the same
shared G1, for twelve actual H proposals. Each admitted descendant is consumed
on the fixed pulse-fast and volatile-reset selection tasks. H0 supplies the
corresponding control in each lineage. Thus selection contains 24 descendant
cells and six H0 cells. The registered ranking prefers more complete-service
passes, then fewer nominal judgment disagreements, then fewer charged requests,
then the lower descendant index. Proposal costs and both selection attempts
enter the request criterion; unselected sibling costs remain in the pipeline.
All 30 effective selection cells pass with no nominal judgment disagreement.
All six choices are frozen before final exposure~\cite{selection}.

Each selected descendant and H0 is then consumed by fresh design episodes on
three frozen final inputs: a stretched-port-window transient task, a half-power
C-lite query derivative, and a new O waveform. The C-lite branch is the
existing single-query digital abstraction, not a new RF or physical environment.
The final designer can inspect the complete current task and acquire its
sources, including the new O waveform. Formation and selection do not see these
final inputs, and final feedback cannot trigger reselection. O0 is excluded
from G/H formation and selection. Array projection is the same in all arms.

Later D episodes have at most 12 provider requests, 40 tool calls, 800 s and
two distinct designs. Each H proposal has an eight-request/400 s ceiling.
The one G1 formation has an eight-request/400 s ceiling; its actual three
requests are charged once and included in G1 pipeline cost. The nine final
cells per arm represent three fixed tasks consumed in three descendant-lineage
contexts, not nine independent task families~\cite{protocol,analysis}.

\section{Observed O0 behavior}
\subsection{Event success is shared with the rich-state control}
\begin{table}[htbp]
\centering\small
\caption{Three actual O0 development episodes. Matched counts concern eight
transitions on one task. Candidate and rich-state share one digital response.}
\label{tab:o0}
\begin{tabular}{@{}lrlccr@{}}\toprule
Arm & Poll (s) & Matched & Host O & Self-judgment & Requests \\\midrule
Candidate & 1 & 8/8 & PASS & UNKNOWN & 8 \\
Fixed periodic & 60 & 0/8 & FAIL & FAIL & 5 \\
Rich-state & 1 & 8/8 & PASS & PASS & 5 \\
\bottomrule\end{tabular}
\end{table}

Candidate and rich-state submit identical designs: 15 mF, one-second sampling,
thresholds 0.05/0.01 g, NVM, one-second retry, zero backoff and the confirmed-event
program. They share the same execution key and response hash. The external
event evaluator passes all eight transitions, energy and receiver checks,
with no missing keepalive or idle false event. The shared response contains
18 accepted packets. Two fresh model submissions of one shared digital
response do not supply two independent realizations~\cite{episodes}.

The periodic control's energy and receiver checks also pass, with 348 accepted
packets. Its messages are periodic snapshots, which provide neither the
required start/stop records nor active-state keepalive coverage. It matches
zero of the eight transitions and fails event service. The contrast shows that
energy feasibility and accepted packet count do not by themselves satisfy this
event contract. Because the fixed rich-state control already passes, O0
establishes no additional event-service benefit for the candidate interface.
The candidate also uses eight provider requests, versus five for each control;
these single-episode counts supply no efficiency advantage or statistical
comparison~\cite{o0}.

\subsection{Retained checked alternative and timing sensitivity}
The actual candidate checks one-second sampling first, then 1.5-second
sampling, and finally submits the one-second design. Both checked responses
are retained in its exposure and cost history. In the 1.5-second check, the
first start record is observed at 4204.070423 s and stamped
4204.139063 s, against the external confirmation time 4203 s. Its timestamp
error is 1.139063 s, exceeding the frozen 1 s tolerance even though
energy and delivery observations alone can appear adequate. This observed
timing violation explains retention of the one-second design; it is a checked
alternative within the same episode, not another model repetition or a later
repair to the accepted design~\cite{episodes}.

\subsection{Participant UNKNOWN and the host event PASS remain distinct}
The candidate's submitted service judgment is UNKNOWN. Its original explanation
identifies unobserved dark intervals after brownout, the absence of an O gap
message to bracket those intervals, inherited A/R first-receive and dark-report
metadata, and unavailable source/hardware calibration. Specifically, the
inherited metadata request reception by 1800 s and dark reporting every 600 s;
the first confirmed O event occurs much later. The participant therefore
declines broader acceptance. Its complete original explanation is preserved
verbatim in the manuscript binding as participant testimony, including its
own numerical interpretations~\cite{first,episodes}.

The registered O evaluator instead asks whether the required start, active
keepalive and stop history was delivered under the nominal energy/receiver
model. It returns PASS, and the receipt retains a nominal predicate-disagreement
flag. The participant's physical and broader-reporting UNKNOWN is not a
physical self-judgment error. Both scopes are retained.
Neither this disagreement nor the event PASS resolves the inherited broader
reporting contract. Earlier A/R dark NO\_GO and FIFO FAIL remain historical
failures; O0 supplies no whole-product acceptance.

The one-second observations include dark-period brownouts near 8766.18 and
19587.64 s. The required active intervals do not overlap the subsequent
unobserved gaps. Successful O0 event coverage therefore does not establish
uninterrupted six-hour sensing or successful detection in every possible dark
window. The next waveform must be evaluated as a separate frozen task rather
than inferred from the O0 final state. Physical qualification remains UNKNOWN
for all three arms.

\section{Final results with every cell retained}
\subsection{Two digital tasks saturate; the new O task separates outcomes}
\begin{table}[htbp]
\centering\small
\caption{Final P/F/I counts: PASS, submitted FAIL and INCOMPLETE. The denominator is six digital cells plus three O cells per arm; all 27 remain included.}
\label{tab:final}
\begin{tabular}{@{}lccc@{}}\toprule
Arm & Two digital tasks & New O waveform & Total \\\midrule
H0 & 6/0/0 & 1/1/1 & 7/1/1 \\
G0 & 6/0/0 & 0/1/2 & 6/1/2 \\
G1 & 6/0/0 & 2/0/1 & 8/0/1 \\
\bottomrule\end{tabular}
\end{table}

Both reserved digital inputs have three PASS per arm, for 18/18 final cells
across those two inputs. All final submitted designs pass their energy gate.
On the new O waveform, H0 has one PASS, one submitted FAIL and one INCOMPLETE;
G0 has zero, one and two; G1 has two, zero and one. The complete 27-cell totals
are H0 7/1/1, G0 6/1/2 and G1 8/0/1. Removing incomplete episodes would change
the question and inflate the apparent completion rate~\cite{details,analysis}.

The final per-lineage pass vectors, in lineage order, are H0 $(2,3,2)$,
G0 $(2,2,2)$ and G1 $(3,2,3)$. G1's 8/9 is a descriptive result from three
descendant-lineage contexts conditional on one shared formation. There is no
independent replication of G1 training, no large task sample, and no promotion
of the default H0 or improver. Strong selection performance on two tasks also
does not establish transfer to a different event-time alignment.

\subsection{Encoded event time is stricter than event matching}
\begin{table}[htbp]
\centering\small
\caption{All nine new-O final cells. Every submitted row has energy PASS, eight matched events and 18 accepted packets. Events P/8 counts individually passing events, not matches. Dash denotes no submission.}
\label{tab:o-final}
\begin{tabular}{@{}lrrrrl@{}}\toprule
Lineage-arm & Poll (s) & Retry (s) & Events P/8 & Max error (s) & Service \\\midrule
L1-H0 & 1 & 1 & 6/8 & 1.103 & FAIL \\
L1-G0 & -- & -- & -- & -- & INCOMPLETE \\
L1-G1 & 0.75 & 0.25 & 8/8 & 0.728 & PASS \\
L2-H0 & 0.75 & 5 & 8/8 & 0.728 & PASS \\
L2-G0 & -- & -- & -- & -- & INCOMPLETE \\
L2-G1 & -- & -- & -- & -- & INCOMPLETE \\
L3-H0 & -- & -- & -- & -- & INCOMPLETE \\
L3-G0 & 1 & 0.25 & 6/8 & 1.103 & FAIL \\
L3-G1 & 0.75 & 1 & 8/8 & 0.728 & PASS \\
\bottomrule\end{tabular}
\end{table}

All five O submissions match eight events, receive 18 accepted packets and
pass energy checks. Only three pass complete event service. These three use
0.75 s sampling; the two one-second submissions each pass six individual event
checks and fail two stops with \path{WRONG_EVENT_TIME}. The three successful
designs include H0 as well as G1. Retry periods also differ, so this is an
observed design association, not a controlled one-factor sampling ablation or
an exclusive G1 mechanism~\cite{details}.

\begin{table}[htbp]
\centering\small
\caption{The two failed stops in each one-second submission. These same times occur in both failed designs, whose retry values and execution keys differ. The error is computed from encoded timestamp, not reception time.}
\label{tab:stamp}
\begin{tabular}{@{}rrrr@{}}\toprule
Confirmation (s) & Reception (s) & Encoded time (s) & Error (s) \\\midrule
12077.625 & 12078.588 & 12078.629 & 1.004 \\
17043.625 & 17044.588 & 17044.728 & 1.103 \\
\bottomrule\end{tabular}
\end{table}

Both failed stop receptions occur less than one second after the external
confirmation instant. Nevertheless, their encoded timestamps exceed the one-second
error bound by 0.004 and 0.103 s. Sampling phase and the modeled clock's timestamp
accumulation therefore have to be checked together. Correct matching, prompt
reception, an empty pending queue and energy feasibility cannot substitute for
the encoded-time predicate. The 0.75 s submissions have maximum encoded-time
error 0.728 s on this waveform. Their observed success is specific to the
frozen digital input and its timing, not a physical robustness guarantee.

The two submitted final O failures each self-report PASS; their nominal
event-service disagreements are retained. These differ from the original O0
candidate's broader UNKNOWN: in the final failures the host identifies a
specific encoded event-time violation within the submitted nominal contract.
Missing submissions supply no service self-judgment at all.

\subsection{Four token-censored episodes remain in the denominator}
The incomplete final episodes are L1-G0, L2-G0, L2-G1 and L3-H0 on O.
Each last SDK assistant response has stop reason \path{length}, with 32,768
output tokens and 32,768 reported reasoning tokens, and no final submission.
This is a response-level metadata observation; it is not a six-hour device
failure, an energy FAIL, or evidence that the whole episode used only that
many tokens. The four episodes are not retried. Their requests and costs
remain charged, and the outcome remains INCOMPLETE rather than a conditional
success or an omitted row~\cite{details,analysis}.

\section{Host cache incident and explicit protocol continuation}
The later selection stage encounters a host representation error: historical
\path{observations.json} is loaded without its separately sealed
\path{trace.csv}, while assessment requires the trajectory. Three selection
attempts raise \path{KeyError: 'trace'}; a fourth attempt is truncated by the
hold. These attempts are apparatus failures or censoring, not negative G0/G1
method outcomes. The frozen hold preserves all four records and their costs:
21 provider requests and USD~0.1318473 in settled
upper-bound estimates~\cite{hold}.

The explicit amendment limits repair to reconstructing the already sealed
historical trace and linking four named one-off recovery attempts. Qualification
passes exact-input historical assessment, preservation of compressed raw and
all five previously valid selection results, and rejection of missing trace,
wrong trace hash, wrong input and wrong world. These checks use existing
responses with zero new provider requests or simulations; they establish a
cache decoding boundary, not method success~\cite{qualification}.

Version 1.1 retains the original methods, tasks, scoring and request/time/dollar
budgets, while
preserving every original attempt. Formal selection accounting uses the original
proposal and two effective selection attempts; incident overhead remains
separately itemized and included in full arm/global pipeline cost. The protocol
retains 73 logical slots with a ceiling of 77 actual episodes after four named
recoveries. Request, summed episode-time and settled upper-bound cost ceilings
remain 860, 54,400 s and USD~100. These are ceilings, not observed final totals.
The amendment preserves the five valid selection responses and twelve H
proposals rather than silently replacing them~\cite{amendment,protocol}.
The four named authorized recovery attempts all close with PASS. They resolve
the four canonical slots without rebuilding G/H or erasing the incident
history~\cite{recovery}. Effective selection consequently has 30 cells, while
its full accounting contains 34 actual attempts.

\section{Complete accounting and interpretation}
\begin{table}[htbp]
\centering\small
\caption{Full pipeline accounting from the authoritative analysis. Selection includes the four original host incidents and four recoveries. USD denotes fixed-rate provider upper-bound estimates; wall is summed episode time.}
\label{tab:cost}
\begin{tabular}{@{}lrrrr@{}}\toprule
Stage & Attempts & Requests & USD & Wall (s) \\\midrule
O0 development & 3 & 18 & 0.1579020 & 240.4 \\
G1 formation & 1 & 3 & 0.0269157 & 38.7 \\
All H proposals & 12 & 43 & 0.4491399 & 862.9 \\
Selection, all attempts & 34 & 173 & 1.2049959 & 2028.0 \\
All final cells & 27 & 130 & 1.3377693 & 2781.6 \\
Total & 77 & 367 & 3.1767228 & 5951.6 \\
\bottomrule\end{tabular}
\end{table}

The full pipeline records 367 requests, all SETTLED, with no new UNKNOWN usage,
USD~3.1767228 in fixed-rate settled upper-bound estimates and 5951.639484 s of
summed episode wall time. This includes the one G1 formation, all twelve H
proposals, eliminated siblings, selection, final episodes, four original host
incidents, four recoveries and token-censored responses. The incident subset
is 21 requests, USD~0.1318473 and 216.654733 s; it is already in the totals
and must not be added again~\cite{analysis}.

Full G0 pipeline cost is 135 requests and USD~1.2585825; full G1 cost is
138 requests and USD~1.2447618, including its formation. H0 selection plus
final consumption costs 76 requests and USD~0.5154765, and O0 costs 18 requests
and USD~0.1579020. These four buckets sum to the global total. G0 and G1 pipeline
costs include both proposed descendants rather than only the selected child.
Provider aliases and registered fixed rates bound an estimate, not a provider
invoice or the full end-to-end research cost. Codex assistance and historical
acquisition remain unpriced.

O0's scope is one development input and an explicitly assumed digital apparatus.
The results do not qualify a photovoltaic device, PMIC transient model,
vibration sensor, physical clock, NVM implementation or RF receiver. They also
do not erase prior A/R reporting failures. A/R NO\_GO, dark FAIL, strict FIFO
FAIL and product NOT\_ACCEPTED remain preserved. The independent O event gate
does not cover every inherited A/R metadata, reporting or dark-observation
requirement. Physical qualification remains UNKNOWN throughout.

The study establishes a finite separation: original O0 can be served by the
same design as a strong fixed control; a different waveform exposes encoded-time
violations despite successful energy and event matching; and procedural
completion can fail through response truncation. It supplies neither evidence
of an interface advantage beyond rich-state nor a stable recursive-improvement
mechanism. The unchanged H0's own successful 0.75 s submission and the single
shared G1 formation constrain that interpretation.

\section{Reproducible manuscript and evidence limits}
\path{render_phase4.py} verifies the authoritative closed handoff and its
finite source list, cross-checks actual and final rows, frozen choices,
accounting and per-event timestamp arithmetic, then generates this fragment.
It also preserves the original O0 participant explanation verbatim in
\path{phase4-binding.json}. It runs no model, scientific simulator or service
evaluator. The source-preservation receipt records all six choices frozen
before final exposure, unchanged historical prefixes and preserved old UNKNOWN
billing records~\cite{preservation}.

The paper-final receipts bind both TeX sources, renderer, shared layout/logo,
compiled PDF, fonts, extracted text and actual page-by-page visual review.
Build and rendering validate presentation and source binding; they do not
replicate the scientific batch, establish hardware acceptance or replace the
coordinator's registered repository verification. No scientific code, task,
result, old manuscript, shared style or logo is changed for this paper.
